\documentclass[11pt,letterpaper]{article}
\usepackage[T1]{fontenc}
\usepackage[utf8]{inputenc}
\usepackage{lmodern}
\usepackage[margin=0.88in,headheight=15pt]{geometry}
\usepackage{amsmath,amssymb,amsthm,mathtools}
\usepackage{mathrsfs}
\usepackage{booktabs,array,longtable}
\usepackage[dvipsnames]{xcolor}
\usepackage{microtype}
\usepackage{enumitem}
\usepackage{needspace}
\usepackage{fancyhdr}
\usepackage{xurl}
\usepackage[colorlinks=true,linkcolor=MidnightBlue,citecolor=MidnightBlue,urlcolor=MidnightBlue]{hyperref}
\usepackage{listings}
\hypersetup{pdftitle={Finite additive coding for an explicit Con(ZFC) equation},
 pdfauthor={Prepared for the ProveIt project with ChatGPT assistance},
 pdfsubject={A ten-integer Diophantine consistency equation and a twelve-positive-variable refinement}}
\pagestyle{fancy}
\fancyhf{}
\fancyhead[L]{\small Finite additive coding for Con(ZFC)}
\fancyhead[R]{\small ProveIt research supplement}
\fancyfoot[C]{\thepage}
\setlength{\parskip}{4pt}
\setlength{\parindent}{0pt}
\setlength{\emergencystretch}{2em}
\setlist{itemsep=2pt,topsep=5pt}
\allowdisplaybreaks[2]
\newcommand{\Z}{\mathbb Z}
\newcommand{\N}{\mathbb N}
\newcommand{\Con}{\operatorname{Con}}
\newcommand{\ZFC}{\mathrm{ZFC}}
\newcommand{\B}{\mathcal B}
\newcommand{\capc}{\mathscr C}
\newcommand{\bitand}{\mathbin{\&}}
\newcommand{\ep}{\varepsilon}
\newtheorem{theorem}{Theorem}[section]
\newtheorem{lemma}[theorem]{Lemma}
\newtheorem{proposition}[theorem]{Proposition}
\newtheorem{corollary}[theorem]{Corollary}
\theoremstyle{remark}
\newtheorem{remark}[theorem]{Remark}
\lstset{basicstyle=\small\ttfamily,columns=fullflexible,breaklines=true,
 frame=single,rulecolor=\color{black!20},backgroundcolor=\color{black!2},
 aboveskip=8pt,belowskip=8pt}

\title{\textbf{Finite additive coding for an explicit\\Con(ZFC) equation}\\[7pt]
\large Ten unrestricted integer unknowns, degree \(69\,554\,088\,184\),\\
and a secondary twelve-positive-variable equation of degree \(938\)}
\author{Prepared for Vladimir Reshetnikov's ProveIt repository\\
\small Mathematical synthesis, explicit circuits, and reproducible audits}
\date{23 October 2026}

\begin{document}
\maketitle
\begin{abstract}
We give a fixed polynomial \(F_{10}\in\Z[f,g,h,k,l,w,x,y,z,m]\) for which
ZFC proves that the absence of integer zeros is equivalent to
\(\Con(\ZFC)\), for the repository's calibrated proof predicate.
The polynomial has exact total degree \(69\,554\,088\,184\) and a
division-free evaluator with \(1\,333\) binary arithmetic operations.
The existing ten-integer form in package 31 has degree
\(126496\cdot5^{31}+27064\) and \(1\,442\) operations. The degree therefore
falls by a factor greater than \(8.46\cdot10^{15}\), while the variable
count is preserved and the evaluator becomes smaller.

The new mathematical step is to encode quartic monomials with a
32-element Bose--Chowla \(B_4\) set instead of radix-five exponent
positions. We prove the complete replacement lemma, including the
off-center coefficient bound, arbitrary-input size bounds, and binary
carry argument. Sun's ten-integer bridge and the repository's final
relation combiner are retained with explicit attribution. A separate
Pell-majorant improvement gives one equation in twelve positive integers
of degree \(938\), or three equations of maximum degree \(929\).

The archive supplies all fixed coefficient payloads, executable
arithmetic circuits, exact finite certificates, proof notes, and source
provenance. These are ordinary mathematical reductions with explicit
inherited dependencies; no new end-to-end Lean or Rocq certificate,
proof of consistency, global minimality, or variable-count record is claimed.
\end{abstract}
\clearpage
\tableofcontents
\newpage

\section{Result, scope, and comparison}
\label{sec:results}

Throughout, \(\N=\{0,1,2,\ldots\}\). The statement \(\Con(\ZFC)\) means
that there is no finite derivation of contradiction in the recursive
ZFC proof system calibrated by ProveIt's K215 endpoint. The endpoint
identifies a blank-input halt of the fixed machine with such a derivation.
This convention fixes the arithmetized consistency statement; it is not
an unspecified truth predicate or an existentially selected theory.

\begin{theorem}[Main fixed instance]
\label{thm:main}
Let \(F_{10}\) be the polynomial defined in
Sections~\ref{sec:fixed}--\ref{sec:finalizer}, equivalently by the delivered
file \texttt{generated/equation10.json} and its fixed bindings. Then
\[
\ZFC\ \vdash\
\left[
 \exists(f,g,h,k,l,w,x,y,z,m)\in\Z^{10}\quad F_{10}=0
 \ \longleftrightarrow\ \neg\Con(\ZFC)
\right].
\]
Consequently ZFC proves
\[
\boxed{\quad
 \left(\forall(f,g,h,k,l,w,x,y,z,m)\in\Z^{10},\ F_{10}\ne0\right)
 \ \longleftrightarrow\ \Con(\ZFC).
\quad}
\]
The total degree is exactly \(69\,554\,088\,184\). Its delivered
arithmetic circuit uses \(484\) additions, \(25\) subtractions, and
\(824\) multiplications: \(1\,333\) operations in total.
\end{theorem}

All letters introduced as abbreviations denote explicit polynomials in
these ten supplied variables. Circuit registers are not additional
existential unknowns. There are no variable exponents, digit-extraction
operations, divisions, radicals, unspecified indices, or free coefficient
parameters in the equation.

The archive also gives a second result, proved in
Section~\ref{sec:positive-majorant}: twelve \emph{strictly positive}
unknowns, either a degree-\(938\) equation with \(122\) operations or a
three-equation system of maximum residual degree \(929\), with \(117\)
operations. Positive and unrestricted integer domains must not be
silently interchanged.

\subsection{The reviewed baseline}

The review is pinned to ProveIt commit
\begin{center}
\texttt{22e75dd6b87b1b199bb08c778d7cf075ae5aae08}.
\end{center}
The relevant directory is
\begin{center}\small
\path{SetTheory/Cardinals/docs/reports/hilbert-tenth-problem/con-zfc-diophantine-forms}.
\end{center}
The comparison includes the source-31 through source-37 refinements
available at that commit, not merely an earlier universal equation.
The decisive comparisons are as follows.

\begin{center}
\small
\begin{tabular}{@{}p{0.30\linewidth}p{0.14\linewidth}p{0.29\linewidth}r@{}}
\toprule
Presentation & Domain & Total degree / maximum degree & Operations\\
\midrule
Source 31, one equation & \(10\) in \(\Z\)
 & \(126496\cdot5^{31}+27064\) & \(1\,442\)\\
New additive code, same finalizer & \(10\) in \(\Z\)
 & \(69\,554\,088\,184\) & \(1\,333\)\\
New code, literal Sun finalizer & \(10\) in \(\Z\)
 & \(1\,323\,756\,220\,880\) & \(1\,383\)\\
\midrule
Source 34, low-degree equation & \(12\) positive & \(1\,022\) & \(124\)\\
New majorant, one equation & \(12\) positive & \(938\) & \(122\)\\
Source 34, associated system & \(12\) positive & \(1\,013\) & \(120\)\\
New majorant, three equations & \(12\) positive & \(929\) & \(117\)\\
\bottomrule
\end{tabular}
\end{center}

The old ten-integer degree is exactly
\[
589\,042\,901\,992\,797\,851\,562\,527\,064.
\]
Its ratio to the new degree is approximately
\(8.4688465822\cdot10^{15}\). The evaluator saves \(109\) operations,
about \(7.56\%\). These improvements hold simultaneously at the same
ten-integer variable count.

Other existing forms retain advantages in other metrics. In particular,
the degree-\(1081\), twelve-positive-variable evaluator takes only
\(108\) operations; the new degree-\(938\) presentation does not dominate
that evaluator. Existing nine-positive-variable and low-degree
many-variable forms remain relevant. The source-31 parent itself is a
system of \(22\) quadratic equations in \(31\) positive variables. There
is no single ordering called \emph{simplest} that makes variable count,
degree, circuit size, coefficient storage, and proof complexity equivalent.
The new ten-integer result is a strict improvement on the reviewed
ten-integer baseline in both reported complexity metrics.

\subsection{What is new and what is inherited}

The general reduction to ten unrestricted integer variables belongs to
Sun~\cite{sun2026}. The fixed consistency machine, its compiler, the
source quartic, and the structural three-radical finalizer belong to the
existing ProveIt construction~\cite{repo}. Bose--Chowla additive sets are
classical~\cite{bosechowla}. The contribution here is their application
to the \emph{entire bounded-degree monomial simplex}, together with a
complete proof that the altered source code meets every downstream
hypothesis. The source reversal target is also tightened to \(4W\),
and the fixed powers are shared explicitly.

The positive-variable improvement is independent of the additive code.
It combines the complete coordinate shift in source 37 with the
translated carrier in source 34, and proves a substantially smaller
separating majorant. It uses the same complete coefficient tuple.

The claimed novelty is relative to the reviewed repository
constructions. We do not infer worldwide priority from an absence of
search results or claim an optimal \(B_4\) set.

\section{A completely fixed quartic source}
\label{sec:fixed}

Use the complete source-31 coefficient tuple associated with the
\(121\times8801\), \(943\)-tile compiler. Its four equation ports are
\[
\beta=\texttt{beta},\quad K_0=\texttt{K},\quad
M_C=\texttt{MC},\quad M_F=\texttt{MF\_source}.
\]
The last port is the \emph{already shifted} mask. Replacing only one
member by a coefficient from a smaller tuple would invalidate the
binding, so no such mixture is made here.

The payloads are nonnegative integers encoded as unsigned little-endian
bytes and then gzip-compressed. Their exact bit lengths are
\[
(1\,064\,921,\ 1\,070\,004,\ 1\,064\,921,\ 1\,064\,922).
\]
They occupy \(533\,099\) uncompressed bytes in total. The included
coefficient builder reconstructs them from the supplied \(943\) Wang
tiles and the \(430\) transitions of the fixed \(215\)-state machine.
The payloads, coefficient recipes, and manifests fix actual integers,
not parameters to be chosen after deciding consistency.

Let \(r_1,\ldots,r_{22}\) be the quadratic residuals listed in
Appendix~\ref{app:parent}, with these four fixed ports. In the ordered
31 source coordinates define
\begin{equation}
P(n_1,\ldots,n_{31})
   =\sum_{j=1}^{22}r_j(n_1+1,\ldots,n_{31}+1)^2.
\label{eq:parent}
\end{equation}
The shift is essential: the parent variables are positive, while the
encoded variables are in \(\N\). The inherited compiler and calibrated
machine endpoint prove
\begin{equation}
\exists n\in\N^{31}\ P(n)=0
\quad\longleftrightarrow\quad
\text{blank K215 halts}
\quad\longleftrightarrow\quad
\neg\Con(\ZFC).
\label{eq:inherited}
\end{equation}
The included \texttt{KEndpoint215.lean} records the final theorem as
\texttt{blank\_halts\_iff\_zfc\_refutation}; its imports and the full
compiler are inherited dependencies, not newly rebuilt formal proofs.

The source polynomial has exact degree four. For example, the
\(\tau^4\) coefficient in the sum of squares is nonzero. Also \(P(0)>0\):
at the all-one parent assignment the fifth residual \(E-X(a-E)\) is
one. This positivity will make the reversed encoding polynomial
positive even at inputs that have not yet passed the digit test.

Write
\[
P(n)=\sum_{|\alpha|\le4}p_\alpha n^\alpha,\qquad
c_\alpha=(4-|\alpha|)!\prod_{i=1}^{31}\alpha_i!\,p_\alpha.
\]
The exact reconstruction produces \(402\) nonzero weighted records.
Their coefficients are small polynomial recipes in the four fixed
ports; after expanding those recipes there are \(675\) symbolic terms.
Let \(L=\max_\alpha|p_\alpha|\). The inherited, replayed coefficient
bound is
\[
|c_\alpha|+1<2^{4\,269\,855}.
\]
Retain the fixed extra constant
\begin{equation}
\boxed{\beta_S=2^{4\,269\,882}}.
\label{eq:betas}
\end{equation}
Since \(48\cdot33^4=56\,924\,208<2^{27}\), it satisfies
\begin{equation}
\beta_S>48\cdot33^4L.
\label{eq:coefbound}
\end{equation}
Its exponent is even. Both this strict coefficient margin and the
congruence \(\beta_S\equiv1\pmod3\) will be used.

The authoritative new binding is \texttt{data/INSTANCE.json}. The
upstream instance file is retained for fixed coefficients and provenance;
its old radix-five exponents and old output paths are superseded.

\section{Replacing positional exponents by a finite additive code}
\label{sec:b4}

\subsection{The exact property needed}

A \(B_4\) set is a set whose sums of four elements, with repetitions
allowed and order ignored, uniquely determine their multisets. We need
\(32\) elements, not \(31\): one is the zero padding coordinate for
monomials of degree below four.

Fix the following ordered set:
\begin{equation}
\begin{split}
(w_0,\ldots,w_{31})={}&(
0,51609,55430,68596,78030,89463,90233,110498,\\
&110508,128615,136498,211819,241519,305379,312706,314312,\\
&369115,417382,429683,435112,450230,462314,469608,490199,\\
&501539,507151,524944,550792,559943,573626,607799,610895).
\end{split}
\label{eq:weights}
\end{equation}
Thus \(W=w_{31}=610\,895\), and \(w_0=0\) is not an unknown.
The delivered checker verifies all
\(\binom{35}{4}=52\,360\) unordered four-element multisets and confirms
that their sums are distinct modulo \(1\,048\,575\). This stronger
modular property implies uniqueness of their integer sums.

\begin{lemma}[Full-simplex injectivity]
\label{lem:injectivity}
For \(|\alpha|\le4\), let
\[
e(\alpha)=\sum_{i=1}^{31}\alpha_iw_i.
\]
Then \(e\) is injective on all such exponent vectors.
\end{lemma}
\begin{proof}
Represent \(\alpha\) by \(\alpha_i\) copies of \(w_i\), padded by
\(4-|\alpha|\) copies of zero. Equal sums imply equal four-element
multisets. In particular, each positive weight has the same
multiplicity, so the exponent vectors agree.
\end{proof}

Checking only collisions among the \(402\) source monomials would be
insufficient. The other factor in coefficient extraction contains
\emph{every} monomial of total degree at most four. The lemma covers the
whole simplex and therefore prevents those missing-source terms from
contaminating the distinguished coefficient.

\subsection{Origin and independent verification of the weights}

The weights come from the classical Bose--Chowla construction
\cite{bosechowla}. Work in
\[
\mathbb F_2[t]/(t^{20}+t^3+1),
\]
and let \(\theta\) be the residue of \(t\). The supplied program checks
that its successive powers enumerate all \(2^{20}-1\) nonzero bit
patterns and return to one. This finite check shows at once that the
quotient is a field and \(\theta\) is primitive; no external table of
primitive polynomials is assumed.

Put \(Q=2^{20}-1=1\,048\,575\). The subfield of order \(32\) is
\[
\{0\}\ \cup\ \{\theta^{33825j}:0\le j<31\}.
\]
For \(a\) in this subfield define \(e_a\in[0,Q)\) by
\(\theta^{e_a}=\theta+a\). If two sums of four such logarithms agree
modulo \(Q\), the corresponding products of four factors \(\theta+a\)
agree. Their monic degree-four polynomials have a difference of degree
at most three that vanishes at \(\theta\). Since \(\theta\) has degree
four over \(\mathbb F_{32}\), that difference is zero. Unique
factorization gives equality of the multisets.

Multiply all logarithms by the unit \(208418\) modulo \(Q\), translate
by \(-97920\), and take the sorted representatives in \([0,Q)\).
The result is \eqref{eq:weights}. The translation removes a cyclic
gap of length \(437680\). This was a finite search for a convenient
set, not a proof that \(610895\) is the smallest possible diameter.
The separate exhaustive multiset check certifies the literal list
independently of the finite-field implementation.

\subsection{The coefficient-isolation identity}

Define the ordinary univariate polynomial
\begin{equation}
D_c(T)=\sum_{|\alpha|\le4}c_\alpha T^{4W-e(\alpha)}.
\label{eq:dcode}
\end{equation}
Every exponent is nonnegative. Equality in \(e(\alpha)\le4W\) causes
no problem, so a reversal target \(4W+1\), or the old \(5W\), is unnecessary.

\begin{lemma}[Quartic isolation and coefficient bound]
\label{lem:isolation}
For arbitrary \(n_1,\ldots,n_{31}\),
\[
[T^{4W}]\left(1+\sum_{i=1}^{31}n_iT^{w_i}\right)^4D_c(T)
  =24P(n).
\]
If all \(n_i\ge0\), every coefficient of this product has absolute
value at most
\[
24L\left(1+\sum_i n_i\right)^4.
\]
\end{lemma}
\begin{proof}
The term indexed by \(\gamma\) in the multinomial factor can contribute
to degree \(4W\) with the term indexed by \(\alpha\) in \(D_c\) only if
\(e(\gamma)=e(\alpha)\). Lemma~\ref{lem:injectivity} gives
\(\gamma=\alpha\). Its multinomial coefficient is
\[
\frac{24}{(4-|\alpha|)!\prod_i\alpha_i!},
\]
which cancels the factorial in \(c_\alpha\), proving the identity.

Different \(\alpha\)'s give different exponents of \(D_c\), and
\(|c_\alpha|\le24L\). For any coefficient of the product, bound each
contributing coefficient of \(D_c\) by \(24L\), and sum all the
nonnegative coefficients of the multinomial factor. Their sum is
\((1+\sum_i n_i)^4\). No factor counting the \(402\) source terms is needed.
\end{proof}

Because \(P(0)>0\), the leading coefficient of \(D_c\) is \(24P(0)>0\).
For every integer \(B\) with \(B-1>24L\),
\[
\left|D_c(B)-24P(0)B^{4W}\right|
 \le24L\sum_{j=0}^{4W-1}B^j<B^{4W}.
\]
In particular \(D_c(B)>0\). This estimate holds before any decoding
assumption on the integer being tested.

\section{The new two-variable source code}
\label{sec:coding}

The only supplied variables in this section are \(f,g\in\Z\). Put
\begin{align}
b&=1+3f,& \B&=\beta_Sb^4,&
\ell_c&=\sum_{i=1}^{31}\B^{w_i},\nonumber\\
N_0&=\B^{W+1},&
N_1&=4\B^{8W+1},&
N&=N_0N_1=4\B^{9W+2},\nonumber\\
M_c&=N_0-1-(b-1)\ell_c,&
T_c&=M_c+(\B-2)\B^{4W}N_0,\nonumber\\
J_c&=2(1+g)^4D_c(\B)+\B\sum_{j=0}^{8W}\B^j,&
S_c&=g+J_cN_0,\nonumber\\
R_c&=(S_c+T_c+1)N+T_c+1,&
X&=(N-1)R_c,&Y&=N^2,\nonumber\\
\capc&=b\B^W.
\label{eq:newsource}
\end{align}
Here and below fixed-index sums and fixed integer powers are ordinary
polynomial notation. In the circuit they are evaluated by binary
additions and multiplications, without division.

\Needspace{8\baselineskip}
\begin{proposition}[Replacement source lemma]
\label{prop:source}
The formulas \eqref{eq:newsource} have these properties.
\begin{enumerate}
\item If \(P\) has a nonnegative integer zero, one can choose \(f\)
arbitrarily large, with \(b\) both a square and a power of two, and
choose \(b\le g<\capc\), so that
\[
Y\mid\binom{2X}{X}.
\]
\item Conversely, \(f\ne0\), \(0\le g<2\capc\), \(b\) both a square
and a power of two, and this divisibility imply that \(P\) has a
nonnegative integer zero.
\item Even before assuming that \(b\) is a power of two, the conditions
\(f\ne0\), \(b\) a square, and \(0\le g<2\capc\) imply
\[
0\le S_c,T_c<N,\qquad 3\mid X,\qquad
X>3b,\qquad Y>\max(b,256).
\]
\end{enumerate}
\end{proposition}

\subsection{Bounds before decoding}

Since \(b=1+3f\) is a square and \(f\ne0\), it is at least four and
\(f>0\). Inequality \eqref{eq:coefbound} gives \(\B>2b\).
The base-\(\B\) digits of \(M_c\) are \(\B-b\) at the selected positions
\(w_i\), and \(\B-1\) at all other positions below \(W+1\).
Consequently
\[
0\le M_c<N_0,\quad 0\le g<N_0,\quad 0\le T_c<N.
\]
The bound on \(g\) uses \(g<2b\B^W<N_0\).

For arbitrary such \(g\), not necessarily a correctly packed witness,
\(g+1\le2b\B^W\). Let \(A_c=48L(2b)^4\). The coefficient margin gives
\(A_c<\B-2\). Positivity and the geometric upper bound on \(D_c\) imply
\[
\begin{split}
0<J_c
&\le A_c\B^{4W}\sum_{j=0}^{4W}\B^j+
       \B\sum_{j=0}^{8W}\B^j\\
&<\frac{(\B-2)+\B}{\B-1}\B^{8W+1}
 =2\B^{8W+1}<N_1.
\end{split}
\]
Thus \(0\le S_c<N\) as well. Because \(\beta_S\equiv b\equiv1\pmod3\),
we have \(\B\equiv N\equiv1\pmod3\) and hence \(3\mid X\).
Finally \(R_c\ge N+1>\B\ge b\), and the displayed sizes of \(N,X,Y\)
give the strict inequalities in the proposition.

The cap retains exactly the form required by the downstream theorem:
\begin{equation}
\capc=2^{4\,269\,882W}b^{4W+1}.
\label{eq:capform}
\end{equation}
Both exponents are fixed positive integers, and the exponent of two
is even.

\subsection{The digit mask and the distinguished coefficient}

Assume now that \(b\) is a power of two. Then \(\B,N_0,N\) are powers
of two. Bitwise conjunction is denoted by \(\bitand\) in the proof
only. The literal digits of \(M_c\) imply
\[
g\bitand M_c=0
\quad\Longleftrightarrow\quad
g=\sum_{i=1}^{31}n_i\B^{w_i}\quad\text{with }0\le n_i<b.
\]
At selected positions the mask allows precisely the low
\(\log_2 b\) bits; elsewhere it allows none. This elementary mask
statement is valid for arbitrary increasing positions, as in
Sun's Lemma 2.4~\cite{sun2021}; see also \cite{bayer2026}.

For such decoded digits put
\[
K_c=
\left(1+\sum_i n_i\B^{w_i}\right)^4D_c(\B)
+\frac{\B}{2}\sum_{j=0}^{8W}\B^j.
\]
By \eqref{eq:coefbound} and Lemma~\ref{lem:isolation}, every coefficient
before padding has absolute value strictly less than \(\B/2\).
Therefore each displayed base-\(\B\) digit of \(K_c\) is an integer
strictly between zero and \(\B\), with no carries or borrows.
Its digit at position \(4W\) is \(\B/2+24P(n)\).
Among the integers strictly between zero and \(\B\), only \(\B/2\)
has zero conjunction with \(\B/2-1\). Thus
\[
K_c\bitand\bigl((\B/2-1)\B^{4W}\bigr)=0
\quad\Longleftrightarrow\quad P(n)=0.
\]

Since \(J_c=2K_c\), shifting both operands by one binary bit gives
\begin{equation}
J_c\bitand\bigl((\B-2)\B^{4W}\bigr)=0
\quad\Longleftrightarrow\quad P(n)=0.
\label{eq:binaryshift}
\end{equation}
It would be incorrect to claim that doubling preserves all base-\(\B\)
digits without carries. Identity \eqref{eq:binaryshift} uses a
\emph{binary shift}, so those possible base-\(\B\) carries do not matter.

Because \(g,M_c<N_0\), the low and high binary blocks separate:
\[
S_c\bitand T_c=0
\quad\Longleftrightarrow\quad
g\bitand M_c=0
\ \text{and}\
J_c\bitand\bigl((\B-2)\B^{4W}\bigr)=0.
\]
The carry/binomial identity of Sun~\cite[Lemma 2.2]{sun2021}, for any
power of two \(N\) and \(0\le S_c,T_c<N\), states that the definition of
\(R_c,X,Y\) in \eqref{eq:newsource} gives
\[
S_c\bitand T_c=0
\quad\Longleftrightarrow\quad
N^2\mid\binom{2X}{X}.
\]
There is no monomial-position assumption in this last identity.
This proves the converse in Proposition~\ref{prop:source}.

For the forward direction, start with a zero \(n\) and choose
\(b=4^r\) arbitrarily large, exceeding all its coordinates. Then
\(f=(b-1)/3\) is a positive integer and can be arbitrarily large.
Pack \(g=\sum_i n_i\B^{w_i}\).
Since \(P(0)>0\), at least one coordinate is nonzero, so \(g\ge\B\ge b\).
The usual finite-digit bound gives \(g<b\B^W=\capc\).
The mask, coefficient test, and carry identity prove the required
divisibility. This completes all parts of the replacement lemma.

\section{The unchanged ten-integer bridge}
\label{sec:bridge}

Introduce the remaining eight supplied integer variables
\(h,k,l,w,x,y,z,m\). With \(X,Y,b,\capc\) as above, define
\begin{align}
L_b&=lY,\quad U=2L_bX,\quad V=4gwY,\quad W_b=bw,\nonumber\\
A&=U(V+1),\quad B_b=2X+1,\quad C=B_b+(A-2)h,\nonumber\\
P_n&=U^2V,\quad K=X+1+k(P_n-2),\nonumber\\
S&=2A-5,\quad T=3W_bC-2(W_b^2-1),\nonumber\\
R_2&=(P_n^2-4)K^2+4,\nonumber\\
\Delta_b&=8\capc^3gK^2
 -g^2\{32(C-KL_b)^2\capc^3+g^2K^2\},\nonumber\\
O&=f^2l^2\Delta_b,\quad \xi=zx,\nonumber\\
D&=(A^2-4)C^2+4,\quad E=C^2D\xi,\nonumber\\
F&=4(A^2-4)E^2+1,\quad
G^*=2(A+2)(A-2)^2E^2-1,\nonumber\\
\Omega&=16E,\quad H^*=C+(B_b-C)F+\Omega Fy,\nonumber\\
I^*&=((G^*)^2-1)(H^*)^2+1.
\label{eq:bridge}
\end{align}
These are Sun's bridge and modified Lucas certificate
\cite[Propositions 6.1--6.2 and Theorems 3.1, 4.1]{sun2026}, in the
source-31 factorization. The local subscript \(b\) merely distinguishes
these symbols from earlier bounds. In particular, \(L_b\) is not the
coefficient bound \(L\), and \(B_b\) is not the coding base \(\B\).

For clarity, the exact imported contract is the following.

\begin{proposition}[Sun bridge contract, specialized]
\label{prop:bridge}
For source polynomials satisfying Proposition~\ref{prop:source} and
cap \eqref{eq:capform}, the formulas \eqref{eq:bridge} have these properties.
\begin{enumerate}
\item A zero of \(P\) yields \(f,g,h,k,l,w\), with \(f\) as large as
needed and \(A>2\), such that \(b\) and \(R_2\) are integer squares,
\(S\mid T\), \(C=u_{B_b}(A)\), and \(O>0\).
\item Given these positive bridge witnesses, the modified Lucas
interpolation theorem permits \(\xi>0\) divisible by any prescribed
positive integer, and permits an arbitrary residue class for \(y\)
modulo \(\xi\), while making \(DFI^*\) a square.
\item Conversely, \(b\) a square, \(O>0\), \(\xi\ne0\),
\(DFI^*\) and \(R_2\) squares, and \(S\mid T\), imply a nonnegative
zero of \(P\).
\end{enumerate}
\end{proposition}

Here \(u_j(A)\) denotes the Lucas sequence with \(u_0=0,u_1=1\), and
\(u_{j+1}=Au_j-u_{j-1}\). It occurs in the proof contract, not as an
operation in the final equation. The reverse bridge first obtains
\(g>0\), \(K\ne0\), and \(0<g<2\capc\) from \(\Delta_b>0\), then
recovers the power-of-two and binomial conditions and invokes the
source lemma. Its numerical hypotheses \(3\mid X\), \(X>3b\), and
\(Y>\max(b,256)\) were re-established above, before decoding.
Thus the use of the bridge does not assume that radix-five positions
are still present.

On forward witnesses, \(\Delta_b\) is a positive integer multiple of
\(g\), and \(g\ge b\ge4\). Hence
\[
\Delta_b\ge4,\qquad O=f^2l^2\Delta_b\ge4.
\]
Also \(S=2A-5\) is an odd, nonzero integer at \emph{every} integer
assignment, without any solution hypotheses.

\section{The final polynomial and proof of equivalence}
\label{sec:finalizer}

For reviewability we restate the source-31 finalizer and its proof.
This section is inherited structure with the new source code substituted.

\subsection{A divisibility and nonvanishing merger}

Let
\[
\pi(v)=(v^3-v-1)(v^2+23),\qquad
\Sigma=\xi S^5,\qquad r=T+Sy,
\]
and define
\begin{equation}
\Theta=(r^3-rS^2-S^3)(r^2+23S^2).
\label{eq:theta}
\end{equation}
The polynomial \(\pi\) has no rational root and has a root modulo
every positive integer. Here is a short proof of the latter fact.
For an odd prime \(p\ne23\), the cubic has nonzero discriminant \(-23\).
If it has a root, that root is simple. Otherwise it is irreducible;
Frobenius acts as an even three-cycle, so its discriminant is a
nonzero square modulo \(p\), and \(v^2+23\) has a simple root.
At \(p=23\), the cubic has the simple root \(3\).
Hensel lifting handles all odd prime powers. Since \(-23\equiv1\pmod8\),
it is a square in the \(2\)-adic integers. The Chinese remainder
theorem completes the argument. The rational-root test rules out a
rational root of the cubic, and the quadratic is positive on \(\mathbb Q\).

Since \(S\ne0\), equation \eqref{eq:theta} never gives \(\Theta=0\).
If \(\Sigma\mid\Theta\), then \(S\mid T\): otherwise select a prime with
\(\alpha=v_p(S)>\beta=v_p(T)\). The two homogeneous factors in
\(\Theta\) have valuations exactly \(3\beta\) and \(2\beta\), by their
unique least-valuation terms. Thus \(v_p(\Theta)=5\beta<5\alpha\),
contradicting \(S^5\mid\Theta\).

Conversely, if \(S\mid T\) and \(\xi\ne0\), write \(T=Sq\), select a
root \(c\) of \(\pi\) modulo \(|\xi|\), and prescribe
\(y\equiv c-q\pmod{\xi}\). Then \(\Sigma\mid\Theta\).
This is precisely the residue freedom retained by
Proposition~\ref{prop:bridge}.

\subsection{Merging two square requirements}

Define three radicands and three weights:
\begin{align}
r_0&=b\{1+(3O-4)b^2z^2\},&
r_1&=R_2,&r_2&=DFI^*,\nonumber\\
v_0&=1,&v_1&=r_0,&v_2&=r_0(P_n^2+K^2+3).
\label{eq:radicands}
\end{align}
If \(z\ne0\) and \(r_0\) is a square, its two displayed factors are
coprime. Therefore \(b\) is a signed square; \(b\equiv1\pmod3\)
excludes zero and negative squares. The possible zero second factor
would force \(|b|=1\), also giving \(b=1\).
Thus \(b\) is a positive square, and
\(1+(3O-4)(bz)^2\) is a square. Since \(bz\ne0\), positivity follows
immediately: if \(O\le0\), this integer is at most
\(1-4(bz)^2<0\). Hence \(O>0\).

Conversely, on forward witnesses \(O\ge4\).
The integer \(d=3O-4>0\) is \(2\pmod3\), hence nonsquare.
A positive Pell unit gives a solution \(u^2-dv^2=1\) with \(v>0\)
divisible by \(b\): raise the unit to its finite order modulo \(b\).
Writing \(v=bz\) gives \(z>0\) and makes \(r_0\) a square.
Prescribe this \(z\) as a divisor of \(\xi\) in the interpolation
theorem and set \(x=\xi/z\). No extra supplied unknown is used.

\subsection{Why the small weights force each square}

Whenever \(z\ne0\), one has \(|r_0|>1\).
If \(r_0=0\) or \(r_0=\pm1\), the defining factors force \(|b|=1\);
then \(b=1,f=0,O=0\) and \(r_0=1-4z^2\le-3\), a contradiction.
Moreover
\[
\sqrt{|r_1|}\le P_n^2+K^2+2.
\]
Indeed \(|r_1|\le(P_n^2+4)K^2+4\), and subtracting this from
\((P_n^2+K^2+2)^2\) leaves
\(P_n^4+K^4+P_n^2K^2+4P_n^2\ge0\).
Consequently
\begin{equation}
|v_1|>\sqrt{|r_0|},\qquad
|v_2|>\sqrt{|r_0|}+|v_1|\sqrt{|r_1|}.
\label{eq:separation}
\end{equation}

If a signed sum \(\sum_{j=0}^2\ep_jv_j\sqrt{r_j}\) is rational,
take an automorphism of the multiquadratic splitting field and the
largest index whose nonzero radical term it changes.
The magnitude of that term is at least \(|v_j|\), while the sum of
all earlier magnitudes is strictly smaller by \eqref{eq:separation}.
Its change cannot cancel. Thus every nonzero radical is fixed by
every automorphism and is rational; a rational square root of an
integer is integral. Zero radicands are already squares.
The argument uses complex absolute values and therefore does not
presuppose that the radicands are nonnegative or independent.

\subsection{The radical-free equation}

Set
\begin{align}
Z_n&=\Sigma m-\Theta,&
A_j&=\Sigma^2v_j^2r_j\quad(0\le j\le2),\nonumber\\
V_n&=Z_n^2+A_2-A_0-A_1,&
W_n&=V_n^2+4Z_n^2A_2-4A_0A_1.
\label{eq:norminputs}
\end{align}
The requested polynomial is
\begin{equation}
\boxed{\qquad F_{10}=W_n^2-16V_n^2Z_n^2A_2.\qquad}
\label{eq:f10}
\end{equation}
Equations \eqref{eq:parent}--\eqref{eq:f10}, the literal weights, and
the supplied fixed integer tuple completely specify it.
The formal identity
\begin{equation}
F_{10}=\prod_{\ep\in\{-1,1\}^3}
\left(Z_n+\sum_{j=0}^2\ep_j\sqrt{A_j}\right)
\label{eq:normidentity}
\end{equation}
is used only in its proof. The evaluator contains no radicals.
An exact expansion of the right side has \(35\) integer polynomial
terms in \(Z_n,A_0,A_1,A_2\); the independent symbolic check compares
it coefficientwise to the left side.

\begin{proof}[Proof of Theorem~\ref{thm:main}, apart from the degree]
Start with \(P(n)=0\). The replacement source lemma and the positive
bridge give the required square and divisibility conditions, with
\(O\ge4\). The Pell construction gives \(z>0\) making \(r_0\) square.
The interpolation theorem supplies \(\xi>0\) divisible by \(z\),
and a \(y\) in the prescribed congruence class, such that \(r_2\) is
square and \(\Sigma\mid\Theta\). Put \(x=\xi/z\).
All three \(r_j\) are now integer squares. Select \(m\) so that
\(\Sigma m-\Theta\) is the negative of one signed sum of their
weighted integer roots multiplied by \(\Sigma\).
One factor of \eqref{eq:normidentity} vanishes, giving \(F_{10}=0\).

Conversely suppose \(F_{10}=0\) at an arbitrary integer assignment.
We always have \(S\ne0\) and \(\Theta\ne0\).
If \(\Sigma=0\), all \(A_j\) vanish and \eqref{eq:f10} becomes
\(\Theta^8\ne0\), impossible. Thus \(\xi=zx\ne0\), in particular
\(z\ne0\). Divide a vanishing factor of \eqref{eq:normidentity}
by \(\Sigma\); it gives a rational weighted sum of the roots of
the \(r_j\). The separation argument makes all of them integer
squares, and the same factor gives \(\Sigma\mid\Theta\).
The merger therefore gives \(S\mid T\), \(b\) square, and \(O>0\).
The reverse Sun bridge and Proposition~\ref{prop:source} give a
nonnegative zero of \(P\).

Composition with \eqref{eq:inherited} proves
\(\exists F_{10}=0\leftrightarrow\neg\Con(\ZFC)\).
All imported results and the displayed arguments are ordinary
mathematics formalizable in ZFC. Classical negation gives the
stated equivalence for lack of solutions. This argument neither
assumes nor establishes which side is true in the standard model.
\end{proof}

\section{Exact degree and the small evaluator}
\label{sec:degree}

\subsection{A symbolic degree calculation}

Let
\[
n=36W+8=\deg N,\qquad c=4W+1=\deg\capc.
\]
All degrees in this section give each of the ten supplied unknowns
degree one and each fixed coefficient degree zero. Propagating degrees
through the displayed formulas gives the following useful checkpoints:
\begin{center}
\small
\begin{tabular}{@{}lr@{\qquad}lr@{}}
\toprule
Expression & Degree & Expression & Degree\\
\midrule
\(D_c(\B)\) & \(16W\) & \(J_c\) & \(32W+4\)\\
\(N,S_c\) & \(n\) & \(R_c\) & \(2n\)\\
\(X\) & \(3n\) & \(Y\) & \(2n\)\\
\(A\) & \(7n+3\) & \(C\) & \(7n+4\)\\
\(P_n\) & \(12n+4\) & \(K\) & \(12n+5\)\\
\(\Delta_b\) & \(28n+3c+14\) &
\(O\) & \(28n+3c+18\)\\
\(r_0\) & \(28n+3c+23\) &
\(r_1\) & \(48n+18\)\\
\(r_2\) & \(616n+340\) &
\(\Sigma\) & \(35n+17\)\\
\(v_2\) & \(52n+3c+33\) &
\(A_2\) & \(790n+6c+440\)\\
\bottomrule
\end{tabular}
\end{center}
These are attained degrees, not merely syntactic upper bounds.
For example, the geometric-sum term uniquely dominates \(J_c\);
the leading term of \(\Delta_b\) is supplied by
\(-32g^2(C-KL_b)^2\capc^3\), with \(KL_b\) dominating \(C\);
and \(K^2\) uniquely dominates the parenthesis in \(v_2\).
Following the nonzero leaders verifies the subsequent rows.

The scaled third radicand \(A_2\) strictly dominates \(A_0,A_1,Z_n^2\).
In the final norm its fourth power has coefficient one and dominates
all other terms. Consequently
\begin{align*}
\deg F_{10}
 &=4(790n+6c+440)\\
 &=3160n+24c+1760\\
 &=113856W+27064\\
 &=\boxed{69\,554\,088\,184}.
\end{align*}
This argument depends on the nonzero positive \(\beta_S\), not on
an accidental cancellation among the four native compiler ports.

As a separate finite attainment certificate, substitute every supplied
unknown by one indeterminate and compute the degree and leading
coefficient modulo three primes. At the claimed degree the leading
coefficients are \(95\pmod{101}\), \(909\pmod{1009}\), and
\(668382\pmod{1000003}\). A single nonzero such coefficient already
proves attainment once the degree upper bound is known. The checker
computes these from the actual fixed tuple.

\subsection{Circuit representation and operation convention}

The JSON circuit is an ordered list of literal, fixed, variable,
addition, subtraction, and multiplication nodes. Every arithmetic
node references earlier nodes. Only the ten listed variable leaves
are quantified. The \(79\) fixed leaves are \(\beta_S\) and the
canonical nonliteral coefficient recipes, bound to the actual
integer payloads. Fixed leaves and reuse are free; each binary
addition, subtraction, and multiplication, including multiplication
by a numeral, is counted. Powers are expanded into multiplication
nodes. The count does not assert that arithmetic on million-bit
coefficients takes unit bit complexity.

For source evaluation, put \(t_0=\B^W\) and
\(t_i=\B^{W-w_i}\). Then
\[
D_c(\B)=
\sum_\alpha c_\alpha t_0^{4-|\alpha|}
                       \prod_{i=1}^{31}t_i^{\alpha_i}.
\]
A finite shared Horner expansion evaluates this polynomial.
The pure powers \(\B^{w_i}\) and \(\B^{W-w_i}\) are constructed
with \(234\) multiplications. The deterministic builder starts with
binary powers, merges frequently shared pairs of disjoint bit
components, and reuses the resulting products. Its certificate
labels each node by its exact exponent; multiplication adds labels,
which proves every target power without relying on numerical tests.

The geometric sum is computed with the division-free identities
\[
G_{2r}=G_r(1+\B^r),\qquad
G_{2r+1}=G_{2r}+\B^{2r},\qquad
G_r=\sum_{j=0}^{r-1}\B^j.
\]
The resulting live arithmetic counts are
\[
484+25+824=\boxed{1333}.
\]
An independent evaluator expands the entire \(D_c\) block exactly
over \(\Z[\B,\beta,K_0,M_C,M_F]\), obtaining the same \(675\)
symbolic terms as the literal definition. It does not import the
circuit builder. It also compares all named formulas on modular
assignments with the actual coefficient tuple.

\subsection{A conservative comparison variant}

The archive retains \texttt{equation10\_sun\_finalizer.json}, which uses
the same new source coding but Sun's literal four-square, sextic
finalizer. Its exact degree is \(1\,323\,756\,220\,880\) and its
evaluator has \(1\,383\) operations. This variant separates the new
additive-coding argument from the repository's later finalizer
optimizations and gives an additional review route.

\section{A secondary improvement: twelve positive unknowns at degree 938}
\label{sec:positive-majorant}

This section uses the complete source-31 coefficient tuple already bundled
with this report, with port aliases
\(\beta=\texttt{beta},K_0=\texttt{K},M_C=\texttt{MC},
M_F=\texttt{MF\_source}\). The last mask is shifted. It does not use the
Sun coefficient or the new exponent code. The inherited normalized
blank-K215 compiler interface is the one proved in source 31's
\texttt{coefficients/COEFFICIENT\_PROOF.md}; source 34's positive reductions
are explicitly uniform in such complete tuples. The new step is a
separator of formal degree nine, replacing the previous degree-23
separator. All symbols in this section are local.

The twelve supplied unknowns are
\[
J,F,\alpha,z,h,i,T,s,w,\eta,\zeta,t\in\Z_{>0}.
\]
All other symbols below are polynomial abbreviations. Set
\begin{align*}
q&=\beta J+1,&\omega&=w+q,&C&=q-F-\alpha,\\
X&=\omega q,&Y&=sq^3,&E&=XY,\\
k&=\eta+\zeta,&c&=kY+\eta,&a&=E+Y,\\
H&=4a+3,&\Delta&=a^2+H,&\ell&=EkY,\\
A&=1+\ell(\ell+k),&G&=1+\Delta c^2,&b_0&=X+ac,\\
I&=ic^2,&D&=1+\Delta I^2,&Q&=(\Delta I)^2,\\
d&=Q-1,&v&=cT,&b&=c+RD.
\end{align*}
Here the abbreviation \(R\), and the two transport expressions, are
\begin{align*}
R&=((q-1)(q-F)+\alpha)(q^2-1)+(M_C+qM_F)J,\\
N_k&=k-hE-R,\qquad
N_t=(K_0+\omega)C+q-F-z(q-1).
\end{align*}
The new majorant and translated carrier are
\begin{gather*}
\boxed{M=k(ik^3+T)^2},\qquad \delta=(a+2)c-\ell,\\
\sigma=b+M\delta+v(a+2)I,\qquad
L=HM+1,\qquad W=b_0+H(t+\sigma).
\end{gather*}
Write \(e=Q(Dv^2+b^2)-1\) and \(B_1=2Qvb\). The polynomial norm is
\begin{equation}
N=\prod_{\substack{\tau=\pm\sqrt A,\;g=\pm\sqrt G\\f=\pm\sqrt D}}
\bigl[d(W-HM\tau-Lg)^2-H^2e+H^2B_1f\bigr].
\label{eq:positive-newnorm}
\end{equation}
This is an integer polynomial; the formal sign product introduces no
quantified unknowns. It is valid also for square, zero, or dependent
radicands. The scalar circuit below and the delivered JSON files give
a literal evaluation using only integer additions and multiplications.
The two presentations are
\begin{equation}
\boxed{N_kN_t(1+N)-1=0},\qquad
\boxed{N_k=1,\quad N_t=1,\quad N=0}.
\label{eq:positive-presentations}
\end{equation}

For completeness, one scalar evaluation of the norm is
\begin{align*}
n&=(HM)^2A,&d_n&=dn,&H_Q&=H^2Q,\\
H_P&=H_QDv^2,&H_U&=H_Qb^2,&H_e&=H_P+H_U-H^2,\\
j&=4H_PH_U,&\mathcal T&=2d(W^2+L^2G),&
\mathcal K&=[d(W^2-L^2G)]^2,\\
a_t&=\mathcal T-2(d_n+H_e),&
b_t&=(d_n-H_e)^2-j-\mathcal K,\\
\bar v&=16d_nj-a_t(a_t\mathcal T+2b_t),&
u&=b_t^2-a_t^2\mathcal K,\\
N&=\mathcal K\bar v^2-u(\mathcal T\bar v-u).
\end{align*}
These are the sourced trace--remainder identities with the new inputs.

\begin{theorem}
For the complete fixed coefficient tuple specified above, ZFC proves
that absence of positive integer solutions to either presentation in
\eqref{eq:positive-presentations} is equivalent to \(\operatorname{Con}(\mathrm{ZFC})\).
The single polynomial has exactly twelve supplied unknowns, degree 938,
and a 122-operation evaluator. The system has three equations, maximum
residual degree 929, and 117 operations. Binary additions, subtractions,
and multiplications by numerals are all charged; reuse and system
comparisons are free.
\end{theorem}

\paragraph{Inherited interface, stated explicitly.}
For the authentic tuple, \(N_k=N_t=\epsilon\),
\(\epsilon\in\{1,-1\}\), implies before recovering any root that
\begin{gather*}
q\ge16,\quad Y\ge q^3,\quad E=XY>R+2,\quad
3q<R<q^4-2,\\
k=R+\epsilon+hE>E,\quad kY<c<k(Y+1),\quad
\sqrt G-\sqrt A>kY^2.
\end{gather*}
The recovered four normalized Pell equations, with integral roots and
\(g\equiv b_0\pmod H\), imply a blank K215 halt; the signed-quotient
extension restores \(\epsilon=1\). Conversely a halt admits canonical
positive roots \(g,\tau,f,y,\gamma\), with \(X=2^R\).
These are the inherited compiler assertions, not claims established by
finite tests in the present report. The coordinate shift is complete:
\(2^R/q>q\), so the original \(w_{\rm old}\) can be replaced by the
positive integer \(w=w_{\rm old}-q\). This is the shift proved in source 37.

\paragraph{The new majorant estimate.}
The shifted coordinate gives \(X\ge q(q+1)>q^2\), hence \(R<X^2\).
The bounds
\(c<(17/16)kY\) and \(a+2<(17/16)E\) follow at once.
Crucially,
\begin{equation}
\frac{R(a+2)^2c^4}{k^8Y^2}<1.
\label{eq:new-majorant-key}
\end{equation}
Indeed the left side equals
\[
\frac R{X^2}\left(\frac{a+2}{E}\right)^2
\left(\frac c{kY}\right)^4\left(\frac Ek\right)^4
\le\left[\frac q{q+1}(1+q^{-2}+2q^{-5})(1+q^{-3})^2\right]^2<1.
\]
For the last inequality, set \(x=q^{-2}\). The two parentheses have
product at most \((1+2x)(1+x)^2<1+5x<1+1/q\), since \(q\ge16\).

Let \(Z=v\sqrt D+b\). Every complex auxiliary root satisfying
\(dy^2=Q(vf-b)^2-1\), \(f^2=D\), obeys the sharp estimate
\begin{equation}
|y|<Z+v+1.
\label{eq:positive-sharp-y}
\end{equation}
To see this, \(\Delta>R+2\), \(I^2\ge c^4\) give
\(Q-b=\Delta I^2(\Delta-R)-c-R>0\), so \(b\le d\), and
\(\sqrt D<\Delta I<d\). Therefore
\[
|y|^2\le\frac{QZ^2+1}d
<\left[Z\left(1+\frac1d\right)\right]^2,
\qquad Z/d<v+1.
\]
The strict middle inequality has numerator \(QZ^2-d>0\).

Since \(D<(a+2)^2I^2\), equation \eqref{eq:new-majorant-key} yields
\(b/(kY^2)<i^2k^7+1\). Also
\[
\frac{v\sqrt D}{kY^2}
<(17/16)^4EYiTk^2
<\frac{(17/16)^4}{X}iTk^4<iTk^4.
\]
Using \(v/(kY^2)<T\) and \(1/(kY^2)<1\), we obtain
\[
\frac{Z+v+1}{kY^2}<i^2k^7+iTk^4+T+2
<i^2k^7+2iTk^4+kT^2=M.
\]
Thus \(|y|<MkY^2<M(\sqrt G-\sqrt A)\).
Moreover \(\sqrt A>\ell=XkY^2\), so
\(M\sqrt A>XZ>2Z\). The latter will ensure positivity in the converse.

\paragraph{Soundness.}
The scalar norm satisfies \(N\equiv\square\pmod4\): write
\(\mathcal T=2x_0\) and
\(\mathcal K=x_0^2-4d^2W^2L^2G\) in the final formula. Thus a zero of
the product in \eqref{eq:positive-presentations} forces \(N=0\) and
\(N_k=N_t=\epsilon\in\{1,-1\}\). A vanishing factor gives coherent
algebraic roots with
\[
x=(W-Lg)/H=M\tau+y,\quad \tau^2=A,\quad g^2=G,\quad f^2=D,
\quad dy^2=Q(vf-b)^2-1.
\]
Every conjugate of \(x\) has magnitude less than \(M\sqrt G\).
If \(g\) were irrational, an automorphism reversing it would produce
two such \(x\)'s at distance \(2L\sqrt G/H>2M\sqrt G\), impossible.
Hence \(g\) is rational and integral, and \(x\) is rational. A sign
change of an irrational \(\tau\) would require two auxiliary roots of
magnitude less than \(M\sqrt A\) to compensate a difference
\(2M\sqrt A\). Thus \(\tau\) is integral, \(y\) is rational, and
\(f=(e-dy^2)/B_1\) is rational and integral.

The sourced rational Pell saturation lemma now applies to
\[
(\Delta I)^2(vf-b)^2-((\Delta I)^2-1)y^2=1
\]
and makes \(y\), hence \(x\), integral. This lemma says that for
integers \(S\ge2,V\), a rational \(y\) satisfying
\(S^2V^2-(S^2-1)y^2=1\) is integral. Source 24 proves it by reducing
the denominator to a common Pell unit and using divisibility of its
cosine sequence.

Since \(\delta,\sigma>0\), \(W>0\); a negative \(g\) would contradict
the bound on \(x\). Reducing \(W-Lg=Hx\) modulo \(H\) gives
\(g\equiv b_0\pmod H\). Its quotient \(\gamma\) is positive because
\(g>(a+1)c>b_0\). The normalized compiler interface now supplies a blank
K215 halt; its ordinary strong coordinate is \(i_{\rm ord}=\Delta i\).
No independent sign changes of dependent radicals were assumed.

\paragraph{Completeness.}
From a canonical positive compiler witness and the complete \(w\)-shift,
put
\[
t=\gamma+M(g-\delta)+M\tau+y-b-v(a+2)I.
\]
We have \(0<\delta<ac<g\), because \(\ell/c>E/2>2\).
Also \((a+2)I<2\sqrt D\) and
\(M\tau>2(v\sqrt D+b)>b+v(a+2)I\). Therefore \(t\) is a positive
integer. Its definition makes \(W=Lg+H(M\tau+y)\), so one norm factor
vanishes and \(N_k=N_t=1\). Composing both directions with the calibrated
K215 refutation theorem, and negating existential solvability, proves the
stated ZFC equivalence.

\paragraph{Exact degree.}
Scale the twelve variables by a formal parameter. The degrees of
\(E,H,c,I,D,Q,b,M,v\) are respectively
\(6,6,5,11,34,46,38,9,6\).
The positive-leading roots \(g,\tau,f\) have degrees \(11,11,17\), with
\[
g-\tau=\delta-k/2+O(\lambda^{-1}),\qquad
(a+2)I-f=I/(2(a+2))+\text{lower-degree terms}.
\]
The four positive-\(f\) factors in \eqref{eq:positive-newnorm} have degree
119. In the negative-\(f\) factors the leading auxiliary terms cancel.
Their remaining leading terms have degrees 116,116,114,107: the first
three come from the main radical signs, and the last from
\(Hv((a+2)I-f)\), of degree 17. All competing squared terms have degree
at most 104. Thus \(\deg N=4(119)+2(116)+114+107=929\).

More explicitly, with \(q_0=\beta J\), \(k_0=\eta+\zeta\),
\(\omega_0=w+q_0\), the nonzero homogeneous leader is
\[
N_{929}=-2^{46}i^{43}k_0^{103}T^5s^{152}
\omega_0^{69}q_0^{549}(q_0-F)^8.
\]
This is nonzero for every authentic \(\beta>0\). Since
\(\deg N_k=7\) and \(\deg N_t=2\), the product has exact degree 938.
The new product improves source 34's degree-1022, 124-operation form in
both metrics; its system improves the corresponding degree-1013,
120-operation system. The existing degree-1081 evaluator at 108 operations
retains its operation-count advantage. There is no variable-count or
global-optimality claim.

The delivered checker compares 123 independent formal-radical products
with the scalar circuit and verifies six complete univariate degree
specializations, including three with the actual coefficient tuple.
The generic runs also check the displayed leading coefficient. The full
proof details are in \texttt{positive/POSITIVE\_PROOF.md}. These finite
checks support the algebra and do not replace the inherited compiler
theorems or claim a new end-to-end formalization.


\section{What the finite artifacts establish}
\label{sec:verification}

The archive distinguishes mathematical proofs, exact finite
certificates, and diagnostic evaluations.

\paragraph{Unbounded mathematical arguments.}
The full-simplex coefficient identity, arbitrary-input source bounds,
carry reduction, and finalizer equivalence are proved above. The
positive-majorant proof is included in this report and in a longer
proof note. These are combined with explicitly identified compiler,
Pell, Lucas, and machine theorems. Those imported arguments are not
replaced by a successful Python run.

\paragraph{Exact finite certificates.}
The primary driver verifies all \(25\) pinned source-31 file hashes and
the \(17\) additional archived proof dependencies,
reconstructs the shifted quartic and all \(402\) weighted records,
checks the actual fixed coefficient bounds and payloads, reconstructs
the finite-field code, exhausts all \(52\,360\) four-element multisets,
proves the shared-power labels, and regenerates the delivered circuits
byte for byte. It checks the three-radical and four-radical norm
identities as exact sparse polynomials. The coefficient compiler
reconstructs all four native payloads from the supplied finite data
and compares their raw bytes. Its replay includes \(69\,230\) basis
checks and all \(943\) tile checks.

\paragraph{Independent algebra and degree checks.}
The separate primary evaluator proves the entire \(D_c\) source-encoding
polynomial identity by sparse expansion, then compares every named
formula at \(60\) assignments over three prime fields. Exact
nonzero leading coefficients certify degree attainment.
The secondary checker performs \(123\) formal-radical product
comparisons and six complete univariate polynomial specializations,
three using the actual coefficient tuple modulo the selected primes.
Its symbolic leading-form proof is given in the report.

\paragraph{What is not certified.}
These scripts do not search the infinite integer domain or establish
whether a zero exists. They do not rebuild the imported Lean
development or mechanize the newly written unbounded proofs.
Independent mathematical review is included as audit notes, but
those notes are not a substitute for a proof-assistant kernel.
The audits were separate AI-assisted reviews during this work, not
external peer review.
The theorem claimed is an ordinary ZFC theorem supported by the
displayed proof and its cited dependencies.

\subsection{Reproduction}

From the extracted top-level directory, with Python 3.10 or later:
\begin{lstlisting}
python3 reproduce.py
\end{lstlisting}
The primary checks use the Python standard library and require no
network access. The optional secondary checker uses NumPy:
\begin{lstlisting}
python3 positive/check_majorant.py
\end{lstlisting}
Build this PDF with a standard LaTeX installation:
\begin{lstlisting}
latexmk -pdf -interaction=nonstopmode -halt-on-error report.tex
\end{lstlisting}
The prebuilt PDF is included. The file \texttt{README.md} gives the
inventory and integration notes. Manifests record the pinned source
paths, Git blob hashes, SHA-256 values, and coefficient conventions.
The upstream payloads and their license notices are preserved.

\section{The general lesson and remaining opportunities}

The old radix-five choice pays for an exponent vector in which every
coordinate could independently range from zero to four. A quartic
only uses vectors with \emph{total} degree at most four. This is the
structural excess removed by the additive code.

More generally, for a degree-\(\delta\) polynomial with \(\delta\ge2\) in \(n\) source
variables, a \(B_\delta\) set of \(n+1\) elements, one translated to
zero, makes
\[
e(\alpha)=\sum_{i=1}^n\alpha_iw_i
\]
injective on \(|\alpha|\le\delta\). With reversal target \(\delta W\)
and coefficients
\((\delta-|\alpha|)!\prod_i\alpha_i!\,p_\alpha\), the same proof
isolates \(\delta!P\). Bose--Chowla supplies such sets in a cyclic
group of order \(q^\delta-1\) for prime powers \(q\ge n+1\).
For fixed \(\delta\), this permits exponent ranges of order
\(n^\delta\), rather than the exponential positional scale
\((\delta+1)^n\).

For this particular quartic, the next concrete mathematical target
would be a smaller certified \(32\)-element \(B_4\) set, or a
support-aware code proven to separate the source support from
\emph{all} competing quartic monomials. Merely checking pairwise
collisions within the support is not sufficient. Any smaller \(W\)
can be substituted into the present proof, with primary degree
\(113856W+27064\), provided the full certificate and source bounds
are retained.

Reducing the number of supplied unrestricted integer unknowns below
ten would require a different contribution. Nothing here achieves
that. Likewise, the best remaining simplification may depend on
whether the intended reader values low degree, a short evaluator,
small coefficients, or a short formal proof. The delivered results
give two concrete improvements with complete definitions, while
keeping those objectives distinct.

\appendix
\section{The literal positive quartic parent}
\label{app:parent}

The parent variables, in the order used for the \(31\) positive weights,
are
\[
\begin{split}
(&J,F,\alpha,z,f,h,i,o,j,s,w,\tau,\eta,\zeta,y,\gamma,\\
 &N,E,P,c,a,X,D,C_2,L,F_2,V,g,I,B_{\rm aux},\Delta).
\end{split}
\]
All are positive integers in the parent. The symbols in this appendix
are local and should not be identified with the ten final variables.
Define the affine abbreviations
\[
q=\beta J+1,\qquad k=\eta+\zeta,\qquad
H=4a+3,\qquad C=q-F-\alpha.
\]
The \(22\) equations, whose left side minus right side gives the
residuals in order, are
\begin{align*}
L&=q^2-1,\\
N&=q(L+1),\\
X&=wq,\\
a&=E+sN,\\
E&=X(a-E),\\
c&=k(a-E)+\eta,\\
P&=E(c-\eta)+k,\\
\tau^2&=P(P-k)+1,\\
D&=X+ac+\gamma H,\\
\Delta&=a^2+H,\\
C_2&=c^2,\\
D^2&=\Delta C_2+1,\\
k&=jc-V+1+hE,\\
jc-V&=(q-F)(N-q)-CL+(M_C+qM_F)J,\\
(K_0+w)C+(q-F)-z(q-1)&=1,\\
F_2&=f^2,\\
V&=of-c,\\
g&=y^2-V^2,\\
I&=iC_2,\\
B_{\rm aux}&=I^2,\\
B_{\rm aux}&=\Delta(F_2-1),\\
y^2-B_{\rm aux}g&=1.
\end{align*}
After substituting the affine abbreviations, each residual is
quadratic. Replacing each listed variable by \(n_i+1\) and summing
the residual squares gives \eqref{eq:parent}. The delivered
\texttt{parent\_quartic31.json} is the machine-readable authority;
the verifier reconstructs the weighted coefficients directly from it.

\section{Artifact map and attribution}

\begin{longtable}{@{}p{0.36\linewidth}p{0.60\linewidth}@{}}
\toprule
Path & Role\\
\midrule
\endhead
\path{report.tex}, \path{report.pdf} & Main definitions, proofs, comparisons,
and reproducibility instructions.\\
\path{data/INSTANCE.json} & Authoritative new fixed instance, parameters,
degree, and paths.\\
\path{data/weights_certificate.json} & Literal weights, finite-field
recipe, and exhaustive multiset certificate.\\
\path{generated/equation10.json} & Primary \(1\,333\)-operation,
ten-integer polynomial.\\
\path{generated/equation10_sun_finalizer.json} & New source code with
the literal Sun finalizer.\\
\path{code/build_equation.py}, \path{code/shared_powers.py} &
Deterministic circuit construction and shared exponent evaluation.\\
\path{code/weights.py} & Independent finite additive-code construction
and exhaustive verification.\\
\path{code/check_equation.py} & Source integrity, coefficient
reconstruction, norm identities, and degree checks.\\
\path{code/check_coefficients.py} & Rebuilds the four complete compiler
coefficients from the supplied tiles and machine.\\
\path{code/check_circuit_independent.py} & Exact sparse source identity
and independent scalar transcription checks.\\
\path{positive/} & Twelve-positive-variable circuit, builder,
checker, binding, and full proof.\\
\path{proofs/} & Full coefficient-coding argument and independent
mathematical audit notes for both improvements.\\
\path{upstream/source31/} & Exact pinned dependencies, complete payloads,
compiler data, and original license notices.\\
\path{upstream/dependencies/} & Additional inherited proof excerpts
and the calibrated K215 endpoint.\\
\bottomrule
\end{longtable}

This supplement was prepared with ChatGPT assistance for incorporation
into ProveIt. The mathematical inputs, authorship of the general
theorems, and inherited source files are credited explicitly.
The new explanatory text and scripts are intended for distribution
under the repository's MIT No Attribution license, with all inherited
third-party notices retained. See \texttt{NOTICE.md}, \texttt{LICENSE},
and the machine-data notices in the archive.

\begin{thebibliography}{9}

\bibitem{bosechowla}
R.~C. Bose and S.~Chowla,
\emph{Theorems in the additive theory of numbers},
Commentarii Mathematici Helvetici \textbf{37} (1962/63), 141--147.
Original report version:
\url{https://www.cs.umd.edu/~gasarch/COURSES/858/S13/BoseChowla.pdf}.

\bibitem{sun2021}
Zhi-Wei Sun,
\emph{Further results on Hilbert's Tenth Problem},
Science China Mathematics \textbf{64} (2021), 281--306.
Explicit source coding, arbitrary-position masks, and carry lemmas:
\url{https://arxiv.org/html/1704.03504v7}.

\bibitem{sun2026}
Zhi-Wei Sun,
\emph{Ten unknowns for Hilbert's tenth problem over the integers},
arXiv:2609.01594v2, 2026 preprint.
In particular Proposition 2.1, Theorems 3.1 and 4.1,
Propositions 6.1--6.2, and Section 7.
\url{https://arxiv.org/html/2609.01594v2}.

\bibitem{bayer2026}
Jonas Bayer, Marco David, Malte Ha{\ss}ler, Yuri Matiyasevich,
and Dierk Schleicher,
\emph{Diophantine equations over \(\Z\): universal bounds and parallel
formalization}, arXiv:2506.20909v3, 22 October 2026.
Definition 1.5 and Lemma 1.6 give another explicit statement of the
arbitrary-position mask property.
\url{https://arxiv.org/html/2506.20909v3}.

\bibitem{repo}
Vladimir Reshetnikov and ProveIt contributors,
\emph{Con(ZFC) Diophantine forms}, repository source collection,
reviewed at commit
\texttt{22e75dd6b87b1b199bb08c778d7cf075ae5aae08}.
Packages 24--31 supply the inherited compiler and Pell interfaces;
packages 31--37 supply the reviewed comparison constructions.
\url{https://github.com/VladimirReshetnikov/ProveIt/tree/22e75dd6b87b1b199bb08c778d7cf075ae5aae08/SetTheory/Cardinals/docs/reports/hilbert-tenth-problem/con-zfc-diophantine-forms}.

\end{thebibliography}
\end{document}

